% GENERATED FILE. Edit canonical lessons, then run npm run generate:books.
% canonical-source-hash: fnv1a64:c45a50f7ee6e90d0
% canonical-lessons: GU-C366-sarat, GU-C366-sakya, GU-C366-sambhavna, GU-C366-atkal, GU-C366-cahvu

\chapter{On One Condition}
\label{ch:gu-on-one-condition}
\hlchaptermodality{366}

\begin{chapteropening}
\textbf{By the end of this chapter:} \emph{I can say what will happen, and on what condition.}

Complete the last lesson of chapter 366: I can say what will happen, and on what condition.
\end{chapteropening}

\section[\texorpdfstring{\gu{શરત} (śarat)}{શરત (śarat)}]{\gu{શરત} (śarat) — a stipulation; a condition}

\label{lesson:GU-C366-sarat}



\begin{quote}
Before the new one: say the Gujarati for actual, then the Gujarati for to link.
\end{quote}

\subsection*{You'll want to know: \gu{શરત}}
\textbf{\gu{શરત}} — \emph{śarat} — \textquotedblleft{}a stipulation; a condition\textquotedblright{}.

\begin{quote}
\gu{જો} \gu{તમે} \gu{મારી} \gu{શરત} \gu{માનશો} \gu{તો} \gu{હું} \gu{આ} \gu{કામ} \gu{કરીશ}.
\end{quote}

\begin{quote}
\emph{jo tame mārī śarat mānśo to hũ ā kām karīś.}
\end{quote}

\begin{quote}
If you accept my condition, I will do this job.
\end{quote}

A feminine noun.

\begin{cousinweb}
\textbf{\gu{શરત}} (\emph{śarat}) is usually traced to Arabic \emph{sharṭ}, \textquotedblleft{}a condition\textquotedblright{}, through Persian.
\end{cousinweb}

\begin{grammarlens}[title={What you've built}]
One more word for saying what happens if.
\end{grammarlens}

\subsection*{Your turn}
\begin{itemize}
\raggedright
  \item \emph{Say it:} \emph{śarat}
  \item \emph{Say it:} \emph{śarat}, once more
  \item \emph{Say it:} \emph{jo tame mārī śarat mānśo to hũ ā kām karīś.} — If you accept my condition, I will do this job.
  \item \emph{Recall:} say the Gujarati for actual, then the Gujarati for to link, then say \emph{śarat} again
\end{itemize}

\begin{tcolorbox}[breakable,colback=teal!4,colframe=teal!35!black,title={Before you move on}]
What does \gu{શરત} mean? (\textquotedblleft{}A stipulation; a condition\textquotedblright{}.) Say it once more.
\end{tcolorbox}
\section[\texorpdfstring{\gu{શક્ય} (śakya)}{શક્ય (śakya)}]{\gu{શક્ય} (śakya) — possible}

\label{lesson:GU-C366-sakya}



\begin{quote}
Before the new one: say the Gujarati for to link, then the Gujarati for a stipulation.
\end{quote}

\subsection*{You'll want to know: \gu{શક્ય}}
\textbf{\gu{શક્ય}} — \emph{śakya} — \textquotedblleft{}possible\textquotedblright{}.

\begin{quote}
\gu{જો} \gu{બધા} \gu{મદદ} \gu{કરશે} \gu{તો} \gu{આ} \gu{કામ} \gu{શક્ય} \gu{છે}.
\end{quote}

\begin{quote}
\emph{jo badhā madad karśe to ā kām śakya che.}
\end{quote}

\begin{quote}
If everyone helps, this job can be done.
\end{quote}

\begin{cousinweb}
\textbf{\gu{શક્ય}} (\emph{śakya}) is a tatsama, borrowed from Sanskrit \emph{śakya}, \textquotedblleft{}able, possible, practicable\textquotedblright{} (Monier-Williams).
\end{cousinweb}

\begin{grammarlens}[title={What you've built}]
One more word for saying what happens if.
\end{grammarlens}

\subsection*{Your turn}
\begin{itemize}
\raggedright
  \item \emph{Say it:} \emph{śakya}
  \item \emph{Say it:} \emph{śakya}, once more
  \item \emph{Say it:} \emph{jo badhā madad karśe to ā kām śakya che.} — If everyone helps, this job can be done.
  \item \emph{Recall:} say the Gujarati for to link, then the Gujarati for a stipulation, then say \emph{śakya} again
\end{itemize}

\begin{tcolorbox}[breakable,colback=teal!4,colframe=teal!35!black,title={Before you move on}]
What does \gu{શક્ય} mean? (\textquotedblleft{}Possible\textquotedblright{}.) Say it once more.
\end{tcolorbox}
\section[\texorpdfstring{\gu{સંભાવના} (sambhāvnā)}{સંભાવના (sambhāvnā)}]{\gu{સંભાવના} (sambhāvnā) — a likelihood}

\label{lesson:GU-C366-sambhavna}



\begin{quote}
Before the new one: say the Gujarati for a stipulation, then the Gujarati for possible.
\end{quote}

\subsection*{You'll want to know: \gu{સંભાવના}}
\textbf{\gu{સંભાવના}} — \emph{sambhāvnā} — \textquotedblleft{}a likelihood\textquotedblright{}.

\begin{quote}
\gu{જો} \gu{વાદળ} \gu{છવાશે} \gu{તો} \gu{વરસાદની} \gu{સંભાવના} \gu{છે}.
\end{quote}

\begin{quote}
\emph{jo vādaḷ chavāśe to varsādnī sambhāvnā che.}
\end{quote}

\begin{quote}
If clouds gather, there is a chance of rain.
\end{quote}

A feminine noun.

\begin{cousinweb}
\textbf{\gu{સંભાવના}} (\emph{sambhāvnā}) is a tatsama, borrowed from Sanskrit \emph{sambhāvanā}, \textquotedblleft{}bringing together; supposing\textquotedblright{} (Monier-Williams).
\end{cousinweb}

\begin{grammarlens}[title={What you've built}]
One more word for saying what happens if.
\end{grammarlens}

\subsection*{Your turn}
\begin{itemize}
\raggedright
  \item \emph{Say it:} \emph{sambhāvnā}
  \item \emph{Say it:} \emph{sambhāvnā}, once more
  \item \emph{Say it:} \emph{jo vādaḷ chavāśe to varsādnī sambhāvnā che.} — If clouds gather, there is a chance of rain.
  \item \emph{Recall:} say the Gujarati for a stipulation, then the Gujarati for possible, then say \emph{sambhāvnā} again
\end{itemize}

\begin{tcolorbox}[breakable,colback=teal!4,colframe=teal!35!black,title={Before you move on}]
What does \gu{સંભાવના} mean? (\textquotedblleft{}A likelihood\textquotedblright{}.) Say it once more.
\end{tcolorbox}
\section[\texorpdfstring{\gu{અટકળ} (aṭkaḷ)}{અટકળ (aṭkaḷ)}]{\gu{અટકળ} (aṭkaḷ) — a conjecture}

\label{lesson:GU-C366-atkal}



\begin{quote}
Before the new one: say the Gujarati for possible, then the Gujarati for a likelihood.
\end{quote}

\subsection*{You'll want to know: \gu{અટકળ}}
\textbf{\gu{અટકળ}} — \emph{aṭkaḷ} — \textquotedblleft{}a conjecture\textquotedblright{}.

\begin{quote}
\gu{જો} \gu{મારી} \gu{અટકળ} \gu{સાચી} \gu{હશે} \gu{તો} \gu{તે} \gu{કાલે} \gu{આવશે}.
\end{quote}

\begin{quote}
\emph{jo mārī aṭkaḷ sācī haśe to te kāle āvśe.}
\end{quote}

\begin{quote}
If my guess is right, he will come tomorrow.
\end{quote}

A feminine noun.

\begin{cousinweb}
\textbf{\gu{અટકળ}} (\emph{aṭkaḷ}) CDIAL lists the Gujarati noun, \textquotedblleft{}conjecture\textquotedblright{}, under a reconstructed \emph{aṭṭakkalā}, \textquotedblleft{}a guess\textquotedblright{} (CDIAL 183).
\end{cousinweb}

\begin{grammarlens}[title={What you've built}]
One more word for saying what happens if.
\end{grammarlens}

\subsection*{Your turn}
\begin{itemize}
\raggedright
  \item \emph{Say it:} \emph{aṭkaḷ}
  \item \emph{Say it:} \emph{aṭkaḷ}, once more
  \item \emph{Say it:} \emph{jo mārī aṭkaḷ sācī haśe to te kāle āvśe.} — If my guess is right, he will come tomorrow.
  \item \emph{Recall:} say the Gujarati for possible, then the Gujarati for a likelihood, then say \emph{aṭkaḷ} again
\end{itemize}

\begin{tcolorbox}[breakable,colback=teal!4,colframe=teal!35!black,title={Before you move on}]
What does \gu{અટકળ} mean? (\textquotedblleft{}A conjecture\textquotedblright{}.) Say it once more.
\end{tcolorbox}
\section[\texorpdfstring{\gu{ચાહવું} (cāhvũ)}{ચાહવું (cāhvũ)}]{\gu{ચાહવું} (cāhvũ) — to want}

\label{lesson:GU-C366-cahvu}



\begin{quote}
Before the new one: say the Gujarati for a likelihood, then the Gujarati for a conjecture.
\end{quote}

\subsection*{You'll want to know: \gu{ચાહવું}}
\textbf{\gu{ચાહવું}} — \emph{cāhvũ} — \textquotedblleft{}to want\textquotedblright{}.

\begin{quote}
\gu{જો} \gu{તમે} \gu{ચાહો} \gu{છો} \gu{તો} \gu{આપણે} \gu{આજે} \gu{નીકળીશું}.
\end{quote}

\begin{quote}
\emph{jo tame cāho cho to āpaṇe āje nīkḷīśũ.}
\end{quote}

\begin{quote}
If you want, we will set off today.
\end{quote}

Past, with \gu{તેણે} (he or she …): \gu{તેણે} \gu{મદદ} \gu{ચાહી}, \emph{teṇe madad cāhī}.

\begin{cousinweb}
\textbf{\gu{ચાહવું}} (\emph{cāhvũ}) CDIAL lists the Gujarati verb, \textquotedblleft{}to long for, to love\textquotedblright{}, under a reconstructed \emph{cāh}, \textquotedblleft{}see, look for, desire\textquotedblright{} (CDIAL 4775); \textquotedblleft{}to want\textquotedblright{} is its everyday sense.
\end{cousinweb}

\begin{grammarlens}[title={What you've built}]
One more word for saying what happens if.
\end{grammarlens}

\subsection*{Your turn}
\begin{itemize}
\raggedright
  \item \emph{Say it:} \emph{cāhvũ}
  \item \emph{Say it:} \emph{cāhvũ}, once more
  \item \emph{Say it:} \emph{jo tame cāho cho to āpaṇe āje nīkḷīśũ.} — If you want, we will set off today.
  \item \emph{Recall:} say the Gujarati for a likelihood, then the Gujarati for a conjecture, then say \emph{cāhvũ} again
\end{itemize}

\begin{tcolorbox}[breakable,colback=teal!4,colframe=teal!35!black,title={Before you move on}]
What does \gu{ચાહવું} mean? (\textquotedblleft{}To want\textquotedblright{}.) Say it once more.
\end{tcolorbox}
